\SourceLeaf{420}%
offenbar
\[
\begin{aligned}
\delta x&=\delta\varrho.\cos(\varphi+90^{0})\\
\delta y&=\delta\varrho.\sin(\varphi+90^{0})
\end{aligned}
\]
oder
\[
\begin{aligned}
\delta x&=-\delta\varrho.\sin\varphi\\
\delta y&=+\delta\varrho.\cos\varphi.
\end{aligned}
\]
Da nun, wenn $\delta\varrho$ unendlich klein ist,
\[
\begin{aligned}
\delta V&=p\delta x+q\delta y\\
&=(-p\sin\varphi+q\cos\varphi)\delta\varrho\\
&=\mp\delta\varrho\BareRoot{(pp+qq)}
\end{aligned}
\]
und $V$ in der Curve selbst $=0$, werden die obern Zeichen gelten, wenn $V$ beim Durchgange durch die Curve von der Linken zur Rechten aus dem positiven ins negative übergeht und umgekehrt; verbindet man diess mit dem, was am Ende des 2. Art. gesagt ist, so folgt, dass die Curve immer nach der Seite convex ist, nach welcher $V$ das Zeichen erhält, welches
\[
Pqq-2Qpq+Rpp
\]
hat.

Ist z. B. die Curve ein Kreis und setzt man
\[
V=xx+yy-aa,
\]
so wird
\[
\begin{gathered}
p=2x,\qquad q=2y,\\
P=2,\qquad Q=0,\qquad R=2,\\
Pqq-2Qpq+Rpp=8yy+8xx=8aa,\\
(pp+qq)^{\frac32}=8a^3,\\
r=\pm a
\end{gathered}
\]
und die Curve nach der Seite convex, wo
\[
xx+yy>aa,
\]
wie gehörig.
